\section{Glossary}
\label{sec:glossary}

\begin{description}[style=nextline,leftmargin=1.6em,itemsep=3pt]
\item[Activation] The act of telling one PowerShell window to use the project's virtual
  environment. Shown by \texttt{(.venv)} at the start of the prompt.
\item[Board] The USC EE\,105 circuit board with its Arduino Nano 33 BLE Sense Lite and OV7675
  camera. In this project it only takes photographs.
\item[Classifier] The part of the pipeline that names a symbol from its picture. The project's
  classifier is a convolutional neural network.
\item[Clone] A complete copy of a repository on your own computer, or the act of making one.
\item[CNN (convolutional neural network)] A neural network designed for images. It applies
  small filters across the picture to detect strokes and corners and combines their responses
  into a decision.
\item[Commit] A recorded state of the repository, with an author, a date and a message; also the
  act of recording one.
\item[Current folder] The folder in which PowerShell carries out a command, shown in the prompt.
\item[Epoch] One pass of the training procedure through all training images.
\item[Frame] One photograph from the camera: $320 \times 240$ pixels, each a gray level from 0
  (black) to 255 (white).
\item[Git] The program that records and exchanges the history of a repository.
\item[GitHub] The website on which the team's repository is stored.
\item[\LaTeX{}] A typesetting system in which a document is written as text with commands, for
  example \code|\frac{a}{b}| for a fraction, and converted to a PDF by a program such as
  \cmd{pdflatex}.
\item[MiKTeX] The distribution of \LaTeX{} for Windows used here.
\item[Package] A unit of software written by others and installed by a package manager:
  \cmd{pip} for Python packages, \cmd{winget} for Windows programs, MiKTeX for \LaTeX{}
  components.
\item[Path] In the narrow sense, the location of a file, such as
  \code|C:\ee105\camera-to-latex\config.py|. In the phrase ``on the path'', the list of folders
  in which Windows looks for a program whose name is typed.
\item[Pipeline] The chain of processing stages from photograph to PDF
  (Figure~\ref{fig:pipeline}).
\item[Prompt] The text that PowerShell prints when it is ready for a command.
\item[Pull, push] To fetch the team's commits from GitHub, and to upload one's own.
\item[Repository] A folder whose history is managed by Git.
\item[Segmentation] The stage that finds the individual symbols in a photograph.
\item[Token] One unit of \LaTeX{} text: a command, a brace or a single character.
\item[Virtual environment] A private installation of Python and its packages inside the
  project folder \file{.venv}.
\end{description}

\section{Summary of commands}
\label{sec:summary}

Table~\ref{tab:cmds} lists the commands of the guide for later reference. All but the first
group are run in the project folder with the environment activated.

\begin{table}[H]
\footnotesize
\noindent\hspace*{-12mm}%
\begin{tabular}{L{11.0cm}L{3.3cm}}
\toprule
\multicolumn{1}{c}{\hd{Command}} & \multicolumn{1}{c}{\hd{Effect}} \\
\midrule
\multicolumn{2}{l}{\emph{Opening sequence, in every new PowerShell window}} \\
\lstinline[style=sum]|cd C:\ee105\camera-to-latex| & Go to the project \\
\lstinline[style=sum]|.\.venv\Scripts\Activate.ps1| & Activate environment \\
\midrule
\multicolumn{2}{l}{\emph{Project}} \\
\lstinline[style=sum]|python -m pytest -q| & Run the automatic tests \\
\lstinline[style=sum]|python -m tools.make_test_lines --render| & Simulated photographs \\
\lstinline[style=sum]|python -m model.make_synthetic| & Training images \\
\lstinline[style=sum]|python -m model.train| & Train the classifier \\
\lstinline[style=sum]|python -m tools.eval_pipeline --images test_set\synthetic --split all| & Score the pipeline \\
\lstinline[style=sum]|python -m tools.review_results| & Review document \\
\lstinline[style=sum]|python -m pipeline.main --image test_set\synthetic\T47.png --accept --pdf| & One photograph to PDF \\
\midrule
\multicolumn{2}{l}{\emph{Git}} \\
\lstinline[style=sum]|git pull| & Fetch team's changes \\
\lstinline[style=sum]|git status| & List changed files \\
\lstinline[style=sum]|git add|, \lstinline[style=sum]|git commit -m|, \lstinline[style=sum]|git push| & Select, record, upload \\
\midrule
\multicolumn{2}{l}{\emph{PowerShell}} \\
\lstinline[style=sum]|cls| & Clear the window \\
\lstinline[style=sum]|explorer .| & Open the current folder \\
\lstinline[style=sum]|exit| & Close the window \\
\bottomrule
\end{tabular}
\normalsize
\caption{The commands used in this guide. Each command is one line.}
\label{tab:cmds}
\end{table}

\section{Removing the installation}

Nothing in this guide alters Windows irreversibly. To remove the project from a laptop, delete
the folder \file{C:\bs ee105} in File Explorer; this removes the code, the environment and all
generated data. To remove the tools, open \menu{Settings}\then\menu{Apps}\then\menu{Installed
apps} and uninstall \emph{Git}, \emph{Python 3.12}, \emph{Python Launcher}, and \emph{MiKTeX}.
Poppler is removed with the command
\code|winget uninstall --id oschwartz10612.Poppler|. The GitHub account is independent of the
laptop and remains.

\section{Sources}

The installation procedure follows the publishers' own instructions, consulted on 24 November
2026.
\begin{itemize}
\item GitHub Docs, \emph{Creating an account on GitHub}:
  \url{https://docs.github.com/en/get-started/start-your-journey/creating-an-account-on-github}
\item GitHub Docs, \emph{Inviting collaborators to a personal repository}:
  \url{https://docs.github.com/en/repositories/managing-your-repositorys-settings-and-features/repository-access-and-collaboration/inviting-collaborators-to-a-personal-repository}
\item GitHub Docs, \emph{Caching your GitHub credentials in Git} (Git Credential Manager):
  \url{https://docs.github.com/en/get-started/git-basics/caching-your-github-credentials-in-git}
\item Microsoft Learn, \emph{Use the WinGet tool to install and manage applications}:
  \url{https://learn.microsoft.com/en-us/windows/package-manager/winget}
\item Python documentation, \emph{Using Python on Windows} (version 3.12):
  \url{https://docs.python.org/3.12/using/windows.html}
\item MiKTeX, \emph{Install MiKTeX on Windows}: \url{https://miktex.org/howto/install-miktex}
\item MiKTeX, \emph{Just enough TeX} (automatic installation of packages):
  \url{https://miktex.org/kb/just-enough-tex}
\item Poppler packaged for Windows: \url{https://github.com/oschwartz10612/poppler-windows}
\end{itemize}
The expected outputs of the project's own commands were taken from runs of the repository on the
project lead's Windows laptop on 23 November 2026, and from runs of the same commands in a Linux
test environment on 24 November 2026.
